%\chapter*{附錄：各種末世論觀點}
%\addcontentsline{toc}{chapter}{附錄：各種末世論觀點}


%\appendix
\renewcommand{\appendixname}{附錄}
\renewcommand{\thechapter}{\arabic{chapter}}
\renewcommand{\thefigure}{\appendixname~\thechapter.\arabic{figure}}  % Appendix 1.1, 1.2, …

\refstepcounter{chapter}
\chapter*{附錄：各種末世論觀點}
\addcontentsline{toc}{chapter}%
  {\appendixname~\thechapter\quad 各種末世論觀點}




啟示錄主要有四種解釋方式。它們最大的分別在於：啟示錄所描述的事，是發生在過去、整段歷史、將來，還是每個時代？

\subsection*{未來派（Futurism）}
未來派的特色是「將來」。它認為第4章以後的異象大多尚未發生，要在基督再來前後按次序實現，例如將來會有一位具體的敵基督、一段大災難時期，以及基督再來後在地上的千年國度。與其他三派相比，未來派最按照字面和時間順序來讀經文：過去派認為這些事已經發生，歷史派認為正在歷史中逐步發生，理想派則不把它們對應到特定事件。未來派的長處是重視預言的應驗，提醒信徒警醒等候主再來；難處是使啟示錄的大部分內容與第一世紀的讀者似乎關係不大。

\subsection*{過去派（Preterism）}
過去派的特色是把啟示錄的重心放在「過去」，即約翰當時的第一世紀。它認為書中的預言大多已經應驗，所指的是羅馬帝國對教會的逼迫，以及主後70年耶路撒冷被毀。與其他三派相比，過去派最重視歷史背景和第一世紀的初代教會。約翰明說這些事「必要快成」（啟示錄1:1），所以它首先是寫給當時受逼迫教會的安慰。它的長處是讓經文與初代教會緊密相連；難處是較難解釋書中描述的全面審判與萬物更新，因此多數過去派仍相信基督將來要再來。

\subsection*{歷史派（Historicism）}
歷史派的特色是把啟示錄看作「一部按次序展開的教會歷史」，從使徒時代一直延伸到基督再來。它不像過去派只看第一世紀，也不像未來派只看未來，而是認為書中的異象對應歷史上一連串具體事件與人物，例如七間教會代表教會歷史的七個時期。宗教改革時期這種讀法最為盛行。它的長處是強調神掌管歷史的每一個階段；難處是各家對應的事件往往不同，而且每一代人都容易把自己的時代當成歷史的高潮，所以今天較少學者採用。

\subsection*{理想派／象徵派（Idealism）}
理想派的特色是不把啟示錄對應到任何「特定」的歷史事件，而是把它讀成描繪「每個時代」都存在的屬靈實況：神與撒但的爭戰、教會受苦作見證、邪惡勢力的引誘，以及基督終必得勝。其他三派都在問「這件事何時發生」，理想派卻問「這段異象揭示了什麼屬靈真理」。它的長處是讓啟示錄對每一代信徒都直接適用；難處是可能忽略書中確實具有的歷史背景與應驗的部分。因此許多學者會把理想派與其他觀點結合，例如先看經文在第一世紀的意義，再看它對歷代教會的啟示，最後指向基督的再來。

四者並非互相排斥：一個異象可以先指向羅馬帝國，同時呈現歷代教會的屬靈爭戰，並指向基督最後的得勝。
